% !TEX program = xelatex
\documentclass{misisarticle}
\grnti{16.31.21}
\asjc{1702}
\oecd{1.2}
\title{Устойчивость тематической кластеризации новостей AG News при сокращении обучающей выборки}
% Добавить подтверждённый email в четвёртый аргумент перед подачей.
\misisauthor{Максим Щетинин}{НИТУ «МИСИС»}{Москва, Россия}{}

\abstract{%
При переобучении кластеризации новостей необходимо понимать, насколько сохраняется состав тематических групп. На корпусе AG News исследованы два источника изменчивости метода KMeans при фиксированном представлении TF-IDF. В первой серии трижды удаляли по 20\% обучающих документов при неизменном параметре случайной инициализации. Во второй серии меняли параметр инициализации на полной выборке. Разбиения одной и той же проверочной выборки сравнивали с базовым по скорректированному индексу Рэнда (ARI). При удалении документов ARI составил 0{,}188--0{,}449, при смене инициализации --- 0{,}182--0{,}190. В базовом запуске один кластер включил 514 из 600 проверочных документов, а в терминах центроидов обнаружены следы HTML-разметки. Оценки описывают чувствительность конкретной конфигурации и задают базовый уровень для будущего сопоставления разных представлений текста.}
\keywords{тематическая кластеризация новостей, AG News, TF-IDF, KMeans, устойчивость разбиения, скорректированный индекс Рэнда}

\begin{document}
\maketitle
\section{Введение}
В новостном архиве тематические группы могут использоваться для навигации и подписок. Если при обновлении модели часть архивных заметок недоступна, полезно измерить, насколько сохраняется группировка одних и тех же проверочных новостей.

Цель работы --- оценить устойчивость классической связки TF-IDF и KMeans при удалении документов и отдельно описать изменчивость при смене параметра случайной инициализации. Полученные результаты задают отправную точку для будущего сравнения с представлениями текстов на основе языковых моделей в рамках ВКР.

\section{Связанные работы}
Lange и соавторы~\cite{lange} предложили оценивать разбиения через воспроизводимость кластерной структуры на разных выборках. Здесь эта идея применяется к возмущению уже заданной обучающей части при одном общем наборе для проверки. Dhillon и Modha~\cite{dhillon} исследовали кластеризацию разреженных векторов документов и интерпретацию групп по характерным словам центроидов. Их сферическая модель использует иную геометрию, тогда как в данной работе применён обычный евклидов KMeans. Zhang и соавторы~\cite{zhang} ввели используемый вариант AG News как набор для классификации текстов. В настоящем эксперименте рубрики нужны для формирования исходной подвыборки, а согласие моделей измеряется по их собственным назначениям.

\section{Постановка задачи и данные}
Пусть $V$ --- фиксированное множество проверочных документов, $f_0$ --- базовая модель, а $f_j$ --- модель после изменения условий обучения. Для каждого условия вычисляется $\operatorname{ARI}(f_0(V), f_j(V))$. Исследовательский вопрос состоит в том, какого масштаба изменения вызывают два возмущения при одинаковом представлении текстов и какие группы меняются сильнее.

Из канонического AG News train взято без возвращения по 750 текстов каждой из четырёх рубрик с seed 42. Текст образован соединением заголовка и краткого описания. Повторы проверены по нормализованному регистру и пробелам, совпадений среди 3000 документов не обнаружено. При разбиении с seed 43 сформированы обучающая (2400 текстов) и проверочная (600) части. Официальная часть test не использовалась. Файл train взят из зеркала, указанного в TorchText, и зафиксирован контрольной суммой в артефактах исследования.

\section{Методы}
TF-IDF обучен один раз на 2400 документах. Использованы униграммы, минимальная частота по документам 2, английские стоп-слова scikit-learn и L2-нормировка. Получено 5853 признака. К проверочным документам применено готовое преобразование. Для обучения использован KMeans с четырьмя кластерами, инициализацией k-means++, алгоритмом Lloyd и десятью рестартами. Базовый \texttt{random\_state} равен 44. Число кластеров задано протоколом, рубрики не использованы для обучения или оценки моделей.

В трёх запусках с seed отбора 101, 102 и 103 без возвращения удаляли ровно 480 обучающих строк. KMeans переобучали на оставшихся 1920 векторах с \texttt{random\_state=44}. Во второй серии полные 2400 векторов использовали с \texttt{random\_state} 201, 202 и 203. Все семь моделей предсказывали кластеры для одних и тех же 600 векторов.

ARI учитывает совпадения отношений принадлежности пар к одной группе с поправкой на случайное согласие~\cite{hubert}. Значение 1 означает совпадение разбиений с точностью до перестановки номеров. Для разбора размеров и миграций отдельно применено сопоставление кластеров один к одному по максимальному суммарному перекрытию. Для выбора пограничных примеров использована разность расстояний до двух ближайших центров.

\section{Результаты}
\begin{table}[htbp]
\caption{Согласие проверочных разбиений с базой}
\label{tab:metrics}
\centering
\begin{tabular}{@{}lrrr@{}}
\toprule
Условие (seed) & Train & ARI & Переходы \\
\midrule
База (44) & 2400 & 1,0000 & 0 \\
Удаление (101) & 1920 & 0,3793 & 128 \\
Удаление (102) & 1920 & 0,4491 & 95 \\
Удаление (103) & 1920 & 0,1884 & 185 \\
Инициализация (201) & 2400 & 0,1816 & 193 \\
Инициализация (202) & 2400 & 0,1865 & 190 \\
Инициализация (203) & 2400 & 0,1900 & 186 \\
\bottomrule
\end{tabular}
\end{table}
При удалении с seed 103 согласие с базой близко к серии со сменой инициализации (табл.~\ref{tab:metrics}). Для двух остальных удалений ARI выше. Последний столбец показывает число сменивших группу документов после сопоставления номеров.

\begin{figure}[htbp]
\centering
\includegraphics[width=\columnwidth]{figures/cluster_sizes.png}
\caption{Размеры кластеров на 600 проверочных новостях. Номера сопоставлены с базой.}
\label{fig:sizes}
\end{figure}

В базе размеры групп составили 514, 27, 26 и 33 (рис.~\ref{fig:sizes}). Кластер 1 сохранил те же 27 проверочных документов во всех запусках. Для удаления с seed 101 сильнее всего изменился состав кластеров 0 и 2: из 0 в 2 перешло 110 новостей. При seed 103 из 0 в 2 перешло 107 новостей и ещё 45 из 0 в 3.

\section{Обсуждение}
У кластера 0 высокие веса получили \texttt{39}, \texttt{ap}, \texttt{new}, \texttt{said}, \texttt{reuters}. У кластера 1 выделились \texttt{gt}, \texttt{lt}, \texttt{fullquote}, у кластера 2 --- \texttt{quot} и \texttt{39}. Эти токены могут отражать особенности форматирования и источника, поэтому устойчивый малый кластер сам по себе не свидетельствует о наличии устойчивой темы. В кластере 3 преобладали \texttt{oil}, \texttt{prices}, \texttt{quarter}, \texttt{percent}. Близкие к его центру новости описывали изменения фондового рынка на фоне нефтяных цен.

Заметка о рынке казначейских облигаций перед отчётом о занятости (индекс 103631) в базовой модели находилась у границы групп 3 и 0 (разность расстояний 0,00087) и при первом удалении перешла из 3 в 0. Заметка об оценке ОПЕК причин роста цен на нефть (14950) перешла из 2 в 3 при том же удалении. В её сюжете экономический контекст связан с сообщениями о беспорядках.

\section{Ограничения}
Использовано одно разбиение одного англоязычного корпуса и по три повтора на условие. ARI оценивает согласие разбиений, семантическое соответствие рубрикам отдельно не измерялось. Фиксированный seed при сокращении выборки задаёт ту же настройку генератора, однако изменившаяся матрица может привести к другим стартовым центрам k-means++. Поэтому первая серия измеряет итоговое влияние удаления при заданном seed. IDF полного train сохраняется. Причинный эффект очистки разметки не измерялся.

\section{Заключение}
На фиксированной проверочной выборке обычный KMeans с TF-IDF чувствителен как к удалению части обучающих новостей, так и к смене инициализации. В исследованной конфигурации все шесть изменений дали ARI ниже 0,45. При сравнении с языковыми моделями следует сохранить проверочную часть и выделить очистку текста в отдельное условие.

\begin{thebibliography}{9}
\bibitem{lange} Lange T., Roth V., Braun M.~L., Buhmann J.~M. Stability-Based Validation of Clustering Solutions. Neural Computation, 16(6), 1299--1323, 2004. \href{https://doi.org/10.1162/089976604773717621}{DOI: 10.1162/089976604773717621}.
\bibitem{dhillon} Dhillon I.~S., Modha D.~S. Concept Decompositions for Large Sparse Text Data Using Clustering. Machine Learning, 42, 143--175, 2001. \href{https://doi.org/10.1023/A:1007612920971}{DOI: 10.1023/A:1007612920971}.
\bibitem{zhang} Zhang X., Zhao J., LeCun Y. Character-level Convolutional Networks for Text Classification. Advances in Neural Information Processing Systems 28, 2015. \href{https://arxiv.org/abs/1509.01626}{arXiv:1509.01626}.
\bibitem{hubert} Hubert L., Arabie P. Comparing Partitions. Journal of Classification, 2, 193--218, 1985. \href{https://doi.org/10.1007/BF01908075}{DOI: 10.1007/BF01908075}.
\end{thebibliography}
\end{document}
